% ============================================================
%  Companion -- overview and the master lattice
% ============================================================
\rsection{What this document is}

\begin{roadmap}[Five words, and the arrows between them]
\raggedright
\emph{Has a limit}, \emph{continuous}, \emph{differentiable},
\emph{integrable}, \emph{convergent}. Every course in analysis uses all five,
and the relations among them are the real content of Rudin's Chapters
3--7 --- but they are never assembled in one place, because each is
introduced where it is needed and the connections are made in passing.

\medskip
This document assembles them. For every ordered pair of properties it
records a verdict: either an implication, with a proof, or a
non-implication, with the function that kills it. Nothing is left
ambiguous, and nothing is asserted without one or the other.

\medskip
\textbf{A structural point that shapes everything.} Four of the five
properties are \emph{local} --- they are properties of a function at a
point. Riemann integrability is not; it is a property of a function on a
whole interval, and asking whether $f$ is integrable \emph{at} $p$ is
meaningless. A single five-by-five table would quietly equivocate between
the two. So the material splits into three blocks:

\smallskip
\textbf{Block A} --- the local properties at a point $p$.

\smallskip
\textbf{Block B} --- those properties promoted to ``at every point of
$[a,b]$,'' together with integrability.

\smallskip
\textbf{Block C} --- what survives a limit: given $f_n \to f$, which
properties of the $f_n$ are inherited by $f$?

\medskip
\textbf{The spine.} Block C is not a new subject. It is Blocks A and B asked
a second time with a limit in front, and the answers change --- which is
precisely why uniform convergence had to be invented.

\medskip
\textbf{Numbering.} Implications are \textbf{T1}--\textbf{T18} and are
proved where they are stated. Counterexamples are \textbf{E1}--\textbf{E12};
each has a dossier in the last section. References of the form
\emph{S2.13} or \emph{Exercise 4.18} point into \emph{The Annotated Rudin};
plain numbers like \emph{4.19} are Rudin's own.
\end{roadmap}

\rsection{The lattice, at a glance}

\begin{center}
\begin{tikzpicture}[x=1mm, y=1mm]
  % --- Block A : local
  \node[rmbox, text width=20mm] (dp) at (0,34)  {differentiable\\at $p$};
  \node[rmbox, text width=20mm] (cp) at (43,34) {continuous\\at $p$};
  \node[rmbox, text width=20mm] (lp) at (86,34) {has a limit\\at $p$};
  \draw[rmarrow] (dp) -- node[lbl, anchor=south, inner sep=2.5pt] {T1}
                          node[mlbl, anchor=north, inner sep=2.5pt]
                          {$\not\Leftarrow$ E2} (cp);
  \draw[rmarrow] (cp) -- node[lbl, anchor=south, inner sep=2.5pt] {T2}
                          node[mlbl, anchor=north, inner sep=2.5pt]
                          {$\not\Leftarrow$ E1} (lp);
  % --- Block B : global
  \node[rmbox, text width=20mm] (dg) at (0,4)  {differentiable\\on $[a,b]$};
  \node[rmbox, text width=20mm] (cg) at (43,4) {continuous\\on $[a,b]$};
  \node[rmbox, text width=20mm] (ig) at (86,4) {Riemann\\integrable};
  \draw[rmarrow] (dg) -- node[lbl, anchor=south, inner sep=2.5pt] {T7}
                          node[mlbl, anchor=north, inner sep=2.5pt]
                          {$\not\Leftarrow$ E2} (cg);
  \draw[rmarrow] (cg) -- node[lbl, anchor=south, inner sep=2.5pt] {T8}
                          node[mlbl, anchor=north, inner sep=2.5pt]
                          {$\not\Leftarrow$ E5, E6} (ig);
  % --- vertical links
  \draw[rmarrow] (dp) -- node[lbl, anchor=west, inner sep=2.5pt]
                          {at every point} (dg);
  \draw[rmarrow] (cp) -- node[lbl, anchor=west, inner sep=2.5pt]
                          {at every point} (cg);
  % --- the deep failure
  \draw[rmarrow, dash pattern=on 3pt off 2pt] (ig) -- (lp);
  \node[mlbl, anchor=west] at (100,19) {$\not\Rightarrow$ E7};
\end{tikzpicture}
\end{center}
\figcap{Blocks A and B. Solid arrows are theorems, proved below; the label
beneath each records the failure of its converse and names the function that
witnesses it. Bounded is not enough for integrability (E7), and integrability
implies nothing pointwise --- an integrable function may be discontinuous on
a dense set (E6).}

\begin{signpost}[Read the diagram as three claims]
\textbf{One:} the top row is a strict chain, and each strictness is cheap ---
one explicit function apiece. \textbf{Two:} the bottom row is \emph{not} the
top row repeated. Integrability is genuinely global, and the arrow into it
comes from compactness, not from anything pointwise. \textbf{Three:} there is
no arrow back from integrability. Integrable functions can be badly
discontinuous, and bounded functions can fail to be integrable altogether.
\end{signpost}
